%=============================================================
% Chapitre 3 — Architecture et pipeline ETL
%=============================================================
\chapter{Architecture et pipeline de traitement}
\label{ch:pipeline}

\section{Vue d'ensemble du pipeline en 8 phases}

Le projet s'appuie sur un pipeline de traitement structuré en \textbf{8 phases séquentielles},
allant de la collecte automatisée via internet à la restitution applicative.
\textbf{MongoDB} constitue le socle de stockage central à chaque étape clé du flux :
les données brutes y sont ingérées dès la collecte, la base préparée y est également
stockée, et le pipeline ML s'y connecte directement pour l'entraînement.

\begin{table}[H]
\centering
\caption{Les 8 phases du pipeline Electio-Analytics}
\label{tab:pipeline}
\begin{tabular}{clll}
\toprule
\textbf{Phase} & \textbf{Nom} & \textbf{Outil} & \textbf{Sortie} \\
\midrule
1 & Collecte brute       & Python / wget         & CSV téléchargés \\
2 & Ingestion brute      & PyMongo               & Collection \texttt{raw\_data} \\
3 & Nettoyage            & Pandas + PyMongo      & Données propres \\
4 & Ingénierie $\Delta$  & Pandas / NumPy        & 27 delta-features \\
5 & Stockage final       & PyMongo               & Collection \texttt{data\_final} \\
6 & EDA                  & Seaborn / Matplotlib  & Graphiques \\
7 & Modélisation ML      & Scikit-learn / XGBoost & Modèles .pkl \\
8 & Restitution          & React 19              & Dashboard web \\
\bottomrule
\end{tabular}
\end{table}

\section{Phase 1 --- Collecte des données brutes}

Les données sont collectées automatiquement depuis les portails ouverts de l'\textbf{INSEE}
(Filosofi, RP) via \texttt{wget} ou la bibliothèque Python \texttt{requests}.
Les fichiers couvrent les millésimes \textbf{2016} et \textbf{2022}
au niveau \textbf{canton} (niveau infra-départemental).
Aucune donnée n'est conservée localement de façon permanente : les fichiers
téléchargés sont immédiatement ingérés en base (phase~2) puis peuvent être supprimés.

\medskip\noindent\textit{Téléchargement automatisé depuis internet (wget / Python) :}
\begin{verbatim}
import os, subprocess

URLS = {
    "emploi_2016":    "https://www.insee.fr/data/Emploi_2016.csv",
    "emploi_2022":    "https://www.insee.fr/data/Emploi_2022.csv",
    "logement_2016":  "https://www.insee.fr/data/Logement_2016.csv",
    "logement_2022":  "https://www.insee.fr/data/Logement_2022.csv",
    "chomage_2016":   "https://www.insee.fr/data/Chomage_2016.csv",
    "chomage_2022":   "https://www.insee.fr/data/Chomage_2022.csv",
}

os.makedirs("raw_data", exist_ok=True)

for name, url in URLS.items():
    dest = f"raw_data/{name}.csv"
    subprocess.run(["wget", "-q", "-O", dest, url], check=True)
    print(f"[OK] {name} -> {dest}")
\end{verbatim}

\section{Phase 2 --- Ingestion brute dans MongoDB}

Une fois téléchargés, les fichiers CSV sont immédiatement ingérés dans
\textbf{MongoDB} (collection \texttt{raw\_data}), qui constitue la
\textbf{source de vérité unique} du pipeline.
Chaque document est enrichi d'un champ \texttt{millesime} (2016 ou 2022)
et d'un champ \texttt{theme} (emploi, logement, chômage\ldots) pour faciliter
les requêtes ultérieures.

\medskip\noindent\textit{Chargement CSV et insertion dans MongoDB :}
\begin{verbatim}
import pandas as pd
from pymongo import MongoClient

client = MongoClient("mongodb://localhost:27017/")
db = client["electio_analytics"]
db["raw_data"].drop()     # réinitialisation avant chaque collecte

for name, _ in URLS.items():
    millesime = 2016 if "2016" in name else 2022
    theme     = name.split("_")[0]      # "emploi", "logement", ...

    df = pd.read_csv(f"raw_data/{name}.csv",
                     sep=";", encoding="utf-8")
    df["millesime"] = millesime
    df["theme"]     = theme

    db["raw_data"].insert_many(df.to_dict(orient="records"))

total = db["raw_data"].count_documents({})
print(f"{total} documents insérés dans raw_data.")
\end{verbatim}

\section{Phases 3--4 --- Préparation des données depuis MongoDB}

Toute la préparation (nettoyage, normalisation, ingénierie des
\texorpdfstring{$\Delta$}{delta}-variables) est effectuée en lisant
\emph{directement} la collection \texttt{raw\_data} de MongoDB.
Cette approche garantit la traçabilité complète du flux et évite toute
dépendance aux fichiers CSV locaux.

\medskip\noindent\textit{Lecture depuis MongoDB, nettoyage et calcul des delta-features :}
\begin{verbatim}
import pandas as pd
from pymongo import MongoClient

client = MongoClient("mongodb://localhost:27017/")
db = client["electio_analytics"]

# Chargement depuis la collection brute
df_2016 = pd.DataFrame(list(
    db["raw_data"].find({"millesime": 2016}, {"_id": 0})
))
df_2022 = pd.DataFrame(list(
    db["raw_data"].find({"millesime": 2022}, {"_id": 0})
))

# Nettoyage : suppression des colonnes trop lacunaires
seuil = 0.80
df_2016.dropna(axis=1, thresh=int(seuil * len(df_2016)), inplace=True)
df_2022.dropna(axis=1, thresh=int(seuil * len(df_2022)), inplace=True)

# Ingénierie des 27 delta-features
epsilon = 1e-8
features_2016 = [c for c in df_2016.columns if c.endswith("_2016")]

for col in features_2016:
    col_22 = col.replace("_2016", "_2022")
    if col_22 in df_2022.columns:
        df_merged[f"delta_{col}"] = (
            (df_2022[col_22] - df_2016[col]) /
            (df_2016[col].abs() + epsilon)
        )
\end{verbatim}

La phase 4 produit ainsi \textbf{27 variables de variation inter-recensements}
($\Delta$-features) selon la formule décrite au chapitre précédent (équation~\ref{eq:delta}).

\section{Phase 5 --- Stockage de la base finale dans MongoDB}

La base de données entièrement préparée est stockée dans la collection
\texttt{data\_final} de MongoDB. Cette collection devient la
\textbf{source exclusive} pour les phases aval : EDA et pipeline ML.

\medskip\noindent\textit{Persistance de la base finale dans MongoDB :}
\begin{verbatim}
db["data_final"].drop()
records = df_merged.to_dict(orient="records")
db["data_final"].insert_many(records)
print(f"Base finale : {db['data_final'].count_documents({})} cantons stockes.")
\end{verbatim}

Les quatre collections MongoDB utilisées par le projet sont les suivantes :

\begin{itemize}
  \item Collection \texttt{raw\_data} : données brutes INSEE (2016 et 2022)
  \item Collection \texttt{data\_final} : features nettoyées + 27 delta-variables
  \item Collection \texttt{predictions} : résultats par département et modèle
  \item Collection \texttt{models\_meta} : hyperparamètres et métriques des modèles
\end{itemize}

\section{Phase 7 --- Pipeline ML depuis MongoDB}

L'entraînement est lancé depuis le notebook
\texttt{Nouvelle\_Aquitaine\_Machine\_Learning.ipynb}.
La première opération consiste à charger la collection \texttt{data\_final}
directement depuis MongoDB, sans passer par aucun fichier intermédiaire.

\medskip\noindent\textit{Chargement de la base ML depuis MongoDB :}
\begin{verbatim}
from pymongo import MongoClient
import pandas as pd

client = MongoClient("mongodb://localhost:27017/")
db = client["electio_analytics"]

df = pd.DataFrame(list(db["data_final"].find({}, {"_id": 0})))
X = df.drop(columns=["label"])
y = df["label"]
\end{verbatim}

Le notebook orchestre ensuite les étapes suivantes :

\begin{enumerate}
  \item Chargement depuis MongoDB (\texttt{data\_final}, 40\,000 cantons $\times$ 171 colonnes)
  \item Split stratifié 80\,\%/20\,\% (32\,000 train / 8\,000 test)
  \item Validation croisée 5-fold sur le jeu d'entraînement
  \item Évaluation finale sur le jeu de test
  \item Génération du rapport de classification et de la matrice de confusion
  \item Export des modèles en \texttt{.pkl} et stockage des prédictions dans MongoDB
\end{enumerate}

\section{Architecture globale}

\begin{center}
\begin{tabular}{ccccc}
\fbox{INSEE (internet)} & $\longrightarrow$ & \fbox{wget / Python} & $\longrightarrow$ & \fbox{MongoDB \texttt{raw\_data}} \\[0.3cm]
                        &                   &                      &                   & $\downarrow$ \\[0.3cm]
\fbox{MongoDB \texttt{data\_final}} & $\longleftarrow$ & \multicolumn{3}{c}{\fbox{Préparation ETL (Pandas + PyMongo)}} \\[0.3cm]
$\downarrow$            &                   &                      &                   & \\[0.3cm]
\fbox{Pipeline ML}      & $\longrightarrow$ & \fbox{Prédictions}   & $\longrightarrow$ & \fbox{React 19} \\
\end{tabular}
\end{center}
